% portada-algebra-lineal-ii.tex ────────────────────────────────────────────
% Portada — Notas de clase, Álgebra Lineal II — Universidad de Guanajuato
% Estilo: mismo lenguaje visual que portada-matematicas-elementales.pdf
%   (fondo oscuro, retícula de puntos, bloque de título apilado, cifra
%   gigante translúcida de fondo).
%
% Gráfica: teorema espectral en R^2 — un operador simétrico A manda la
% circunferencia unitaria (dominio) a una elipse cuyos semiejes son las
% direcciones propias de A, con longitudes |λ_1|, |λ_2|. Es una lectura
% geométrica propia de un segundo curso (diagonalización ortogonal,
% formas cuadráticas), más allá del cálculo puntual de Av = λv.
% ────────────────────────────────────────────────────────────────────────────

\documentclass[tikz, border=0pt]{standalone}

\usepackage[T1]{fontenc}
\usepackage[utf8]{inputenc}
\usepackage[spanish]{babel}
\usepackage{ebgaramond}
\usepackage[default]{sourcesanspro}
\usepackage{amsmath}
\usepackage{xcolor}
\usetikzlibrary{calc,arrows.meta}

% ── Paleta (misma familia tonal que las demás portadas de "notas de clase") ──
\definecolor{bg}{HTML}{0B0E1A}
\definecolor{panel}{HTML}{141A2E}
\definecolor{grid}{HTML}{FFFFFF}
\definecolor{teal}{HTML}{2DD4BF}
\definecolor{purple}{HTML}{8B7CF6}
\definecolor{gold}{HTML}{F2C14E}
\definecolor{ink}{HTML}{F4F4F5}
\definecolor{muted}{HTML}{8B93A6}

\begin{document}
\begin{tikzpicture}[font=\sffamily]

  \useasboundingbox (0,0) rectangle (210mm,297mm);

  %% ── Fondo ──────────────────────────────────────────────────────────────
  \fill[bg] (0,0) rectangle (210mm,297mm);

  %% ── Resplandor ambiental detrás del diagrama ──────────────────────────
  \foreach \r/\op in {70/0.05, 52/0.06, 34/0.05} {
    \fill[teal, opacity=\op] (98mm,205mm) circle (\r mm);
  }
  \foreach \r/\op in {55/0.04, 36/0.05} {
    \fill[purple, opacity=\op] (128mm,190mm) circle (\r mm);
  }

  %% ── Retícula de puntos ─────────────────────────────────────────────────
  \foreach \x in {8,15,...,203} {
    \foreach \y in {8,15,...,290} {
      \fill[grid, opacity=0.05] (\x mm,\y mm) circle (0.18mm);
    }
  }

  %% ── Cifra de volumen, gigante y translúcida ────────────────────────────
  \node[font=\fontsize{260}{260}\selectfont\rmfamily\bfseries,
        color=panel, anchor=east] at (200mm,88mm) {II};

  %% ── Barra de identificación superior ────────────────────────────────────
  \fill[purple] (14mm,278mm) rectangle (26mm,279mm);
  \fill[teal]   (27mm,278mm) rectangle (35mm,279mm);

  \node[anchor=east, font=\fontsize{8}{8}\selectfont\sffamily,
        color=muted] at (196mm,278.5mm) {NOTAS DE CLASE~\textperiodcentered~VOL.~II};

  %% ── Diagrama: teorema espectral en R^2 ─────────────────────────────────
  \coordinate (C) at (105mm,198mm);

  % retícula polar tenue
  \foreach \r in {14,28,42} {
    \draw[grid, opacity=0.09, line width=0.3pt] (C) circle (\r mm);
  }
  % ejes tenues
  \draw[grid, opacity=0.12, line width=0.3pt] ($(C)+(-58mm,0)$) -- ($(C)+(58mm,0)$);
  \draw[grid, opacity=0.12, line width=0.3pt] ($(C)+(0,-58mm)$) -- ($(C)+(0,58mm)$);
  \draw[grid, opacity=0.07, line width=0.3pt, dashed] ($(C)+(-41mm,-41mm)$) -- ($(C)+(41mm,41mm)$);
  \draw[grid, opacity=0.07, line width=0.3pt, dashed] ($(C)+(-41mm,41mm)$) -- ($(C)+(41mm,-41mm)$);

  % ecuación
  \node[font=\fontsize{15}{15}\selectfont, color=muted]
    at ($(C)+(0,68mm)$) {$A\mathbf{v} = \lambda\mathbf{v}$};

  % circunferencia unitaria (dominio) — discontinua
  \draw[muted, opacity=0.55, line width=0.6pt, dashed] (C) circle (28mm);

  % elipse imagen, ejes principales rotados 28°
  \draw[teal, line width=1pt, rotate around={28:(C)}]
    ($(C)$) ellipse (44mm and 19mm);

  % direcciones propias (semiejes)
  \draw[gold, line width=1pt, -{Latex[length=2.2mm,width=1.6mm]},
        rotate around={28:(C)}] (C) -- ($(C)+(44mm,0)$)
    node[pos=1, right=1.5mm, font=\fontsize{11}{11}\selectfont, color=gold]
    {$\lambda_1$};
  \draw[purple, line width=1pt, -{Latex[length=2.2mm,width=1.6mm]},
        rotate around={28:(C)}] (C) -- ($(C)+(0,19mm)$)
    node[pos=1, above=1mm, font=\fontsize{11}{11}\selectfont, color=purple]
    {$\lambda_2$};

  \fill[ink] (C) circle (0.6mm);

  %% ── Bloque de título ─────────────────────────────────────────────────
  \fill[gold] (20mm,125mm) rectangle (32mm,126mm);

  \node[anchor=west, font=\fontsize{40}{44}\selectfont\rmfamily\bfseries,
        color=ink] at (19.3mm,109mm) {\'Algebra};
  \node[anchor=west, font=\fontsize{40}{44}\selectfont\rmfamily\bfseries,
        color=teal] at (19.3mm,90mm) {Lineal};

  \node[anchor=west, font=\fontsize{11}{11}\selectfont\rmfamily\itshape,
        color=muted] at (20mm,80mm) {Notas de clase para un segundo curso};

  \draw[grid, opacity=0.12, line width=0.3pt] (20mm,70mm) -- (190mm,70mm);

  \node[anchor=west, font=\fontsize{7.5}{7.5}\selectfont, color=muted]
    at (20mm,63.5mm) {ELABORADAS POR};
  \node[anchor=west, font=\fontsize{13}{13}\selectfont\rmfamily,
        color=ink] at (20mm,56mm) {Ricardo Le\'on Mart\'inez};

  \draw[grid, opacity=0.09, line width=0.3pt] (20mm,48mm) -- (190mm,48mm);

  \node[anchor=west, font=\fontsize{7.5}{7.5}\selectfont, color=muted]
    at (20mm,41.5mm) {PROFESOR};
  \node[anchor=west, font=\fontsize{13}{13}\selectfont\rmfamily\itshape,
        color=ink] at (20mm,34mm) {Herbert Kanarek Blando};

  %% ── Pie de página ────────────────────────────────────────────────────
  \fill[teal] (0,0) rectangle (2.4mm,18mm);

  \node[anchor=west, font=\fontsize{7.5}{9}\selectfont, color=muted]
    at (9mm,11mm) {Departamento de Matem\'aticas};
  \node[anchor=west, font=\fontsize{7.5}{9}\selectfont, color=muted]
    at (9mm,7mm) {Divisi\'on de Ciencias Naturales y Exactas~\textperiodcentered~Universidad de Guanajuato};

  \node[anchor=east, font=\fontsize{15}{15}\selectfont\rmfamily,
        color=gold] at (196mm,8.5mm) {2026};

\end{tikzpicture}
\end{document}
